We provide details on how~\eqref{eq:reflection condition_out_field_generic_vacuum} can be computed.
First we introduce the projector of positive frequency and negative frequency modes for the NS fermion as
\begin{eqnarray}
  P^{(+,\, 0)}(z,w)
  & = &
  \frac{+1}{z-w},
  \qquad
  \abs{z} > \abs{w}
  \\
  P^{(-,\, 0)}(z,w)
  & = &
  \frac{-1}{z-w},
  \qquad
  \abs{z} < \abs{w},
\end{eqnarray}
such that
\begin{equation}
  \oint\limits_{\abs{z} > \abs{w}} \ddw
  P^{(+,\, 0)}(z,w)\,
  \Psi^{(0)}( 0 )
  =
  \Psi^{(0,\, +)}( z ),
\end{equation}
and similarly for the negative frequency modes.
Likewise we introduce the projectors for the field with defects as
\begin{eqnarray}
  P^{(+)}(z,\, w)
  & = &
  \frac{P\qty(z;\, \qty{x_{(t)},\, \rE_{(t)}} )\,
        P\qty(w;\, \qty{x_{(t)},\, -\rE_{(t)}} )
       }{z-w},
  \qquad
  \abs{z} > \abs{w}
  \\
  P^{(-)}(z,\, w)
  & = &
  -
  \frac{P\qty(z;\, \qty{x_{(t)},\, \rE_{(t)}} )\,
        P\qty(w;\, \qty{x_{(t)}, -\rE_{(t)}} )
       }{z-w},
  \qquad
  \abs{z} < \abs{w},
\end{eqnarray}
with $P\qty(z;\, \qty{x_{(t)},\, \rE_{(t)}} ) = \finiteprod{t}{1}{N} \qty( 1- \frac{z}{x_{(t)}} )^{\rE_{(t)}}$ as in the main text.
    
We then compute
\begin{equation}
  \begin{split}
    \qty(P^{(+)}\, P^{(+,\,0)})(z,\, w)
    & =
    \oint\limits_{\abs{z} > \abs{\zeta} > \abs{w}} \ddz
    P^{(+)}(z,\, \zeta)\,
    P^{(+,\, 0)}(\zeta,\, w)
    =
    P^{(+,\, 0)}(z,\, w)
    \\
    \qty(P^{(+)}\, P^{(-,\, 0)})(z,\, w)
    & =
    \frac{P\qty(z;\, \qty{x_{(t)},\, \rE_{(t)}} )\,
          P\qty(w;\, \qty{x_{(t)},\, -\rE_{(t)}} ) -1
         }{z-w}.
  \end{split}
\end{equation}
The last equation is valid when $\rM=\finitesum{t}{1}{N} \rE_{(t)} \le 0$ and for $\abs{z}$ and $\abs{w}$ arbitrary.
Specializing the previous expressions to $\Psi^{(out)}( z )$, we need to constrain $\abs{z} > x_{(1)}$ and $\abs{w} > x_{(1)}$.

Finally the vacuum in presence of defects can be  described by
\begin{equation}
  \begin{split}
    \Psi^{(+)}( z ) \Gexcvacket
    & =
    \qty(P^{(+)}\, \Psi)( z ) \Gexcvacket
    \\
    & =
    \qty(P^{(+)}\, \Psi^{(out)})( z ) \Gexcvacket
    \\
    & =
    \left\lbrace
      \qty(P^{(+)}\, P^{(+,\, 0)}\, \Psi^{(out)})( z )
    \right.
    \\
    & +
    \left.
      \qty(P^{(+)}\, P^{(-,\, 0)}\, \Psi^{(out)})( z )
    \right\rbrace
    \Gexcvacket
    \\
    & =
    0,
  \end{split}
\end{equation}
where we assumed $\abs{z} > x_{(1)}$.
The expression finally becomes~\eqref{eq:reflection condition_out_field_generic_vacuum}.

% vim: ft=tex
